% Generated draft: requires mathematical and source review.
\documentclass[11pt]{article}
\usepackage[margin=25mm]{geometry}
\usepackage{fontspec}
\setmainfont{Latin Modern Roman}
\usepackage{amsmath,amssymb,mathtools,unicode-math}
\setmathfont{Latin Modern Math}
\usepackage{enumitem,needspace,xurl,hyperref}
\setlength{\parindent}{0pt}
\setlength{\parskip}{5pt}
\setlength{\emergencystretch}{3em}
\newcommand{\N}{\mathbb{N}}
\newcommand{\Z}{\mathbb{Z}}
\newcommand{\Q}{\mathbb{Q}}
\newcommand{\R}{\mathbb{R}}
\newcommand{\C}{\mathbb{C}}
\newcommand{\F}{\mathbb{F}}
\newcommand{\A}{\mathbb{A}}
\newcommand{\PP}{\mathbb{P}}
\newcommand{\E}{\mathbb{E}}
\newcommand{\eps}{\varepsilon}
\newcommand{\dd}{\,\mathrm{d}}
\begin{document}
\section*{UnsolvedMath ID: 30005512}
\textbf{Source problem code:} catalogue entry
\subsection{Definitions and mathematical statement}
Let \(p\) be a prime and let \(F\subseteq\C\) be an abelian number field, meaning a finite Galois extension of \(\Q\) with abelian Galois group. For each positive integer \(r\), put \(\Q_r=\Q(\zeta_r)\), where \(\zeta_r=e^{2\pi i/r}\). The conductor of \(F\) is the least positive integer \(n\) such that \(F\subseteq\Q_n\). Write
\[
n=p^a m,\qquad a\in\Z_{\geq0},\quad m\in\Z_{\geq1},\quad \gcd(m,p)=1.
\]
Let \(F\Q_m=\langle F,\Q_m\rangle\) denote the compositum inside \(\Q_n\).

For a finite group \(G\), let \(\operatorname{Irr}(G)\) denote its ordinary irreducible complex characters. For \(\chi\in\operatorname{Irr}(G)\), define its field of values by
\[
\Q(\chi)=\Q\bigl(\chi(g):g\in G\bigr).
\]
Let \(B\) be a \(p\)-block of \(G\), and let \(D\) be a defect group of \(B\). If \(|G|_p=p^u\) and \(|D|=p^d\), where \(t_p\) denotes the largest power of \(p\) dividing a positive integer \(t\), the height \(h_B(\chi)\) of \(\chi\in\operatorname{Irr}(B)\) is determined by
\[
\chi(1)_p=p^{u-d+h_B(\chi)}.
\]
Thus height zero means \(\chi(1)_p=p^{u-d}\).

\textbf{Conjectural statement (supplied as Conjecture B).} For every prime \(p\) and every abelian number field \(F\) with conductor decomposed as above,
\[
\begin{aligned}
&\exists\,\text{a finite group }G,\ \exists\,\text{a }p\text{-block }B\text{ of }G,\ 
\exists\,\chi\in\operatorname{Irr}(B):\\[-2pt]
&\hspace{30mm} h_B(\chi)=0\quad\text{and}\quad\Q(\chi)=F
\end{aligned}
\quad\Longleftrightarrow\quad
\gcd\!\left([\Q_n:F\Q_m],p\right)=1.
\]
The group, block, and character may depend on \(p\) and \(F\).

\textbf{Supplied qualifications.} The original report excerpt begins mid-sentence; the complete hypotheses and quantifiers above follow the supplied reconstructed statement. The excerpt reports that the coprimality condition implies existence of the character, and that the reverse implication is reduced to height-zero characters in blocks of quasisimple groups. Here quasisimple means perfect with quotient by its center a nonabelian simple group. It also reports that the Alperin--McKay--Navarro conjecture implies Conjecture B. These are attributed source claims, not independently verified conclusions of this draft. The supplied dated assessment labels the problem partially solved and reports no general proof. This statement-only formulation requires human mathematical review.
\subsection{Short English statement}
The conjecture characterizes which abelian number fields occur as fields of values of height-zero irreducible characters for a fixed prime \(p\). The proposed criterion is that adjoining the prime-to-\(p\) cyclotomic field leaves an index coprime to \(p\) in the cyclotomic field of the conductor. The supplied report says the existence direction is established and the reverse direction is reduced to quasisimple groups; these claims require independent verification.
\subsection{Sources}
\begin{itemize}
\item Local UnsolvedMath registry, supplied record with permanent ID \(30005512\), title ``Abelian Fields Generated by Height-Zero Characters,'' exact source problem code \texttt{catalogue entry}, in the Oberwolfach Reports Open Problems collection. The record gives proposed year 2023 and names no individual proposer.
\item Supplied report citation: \emph{Representations of Finite Groups} (2023), \emph{Oberwolfach Reports}, pp.~1031--1084, DOI \texttt{10.4171/owr/2023/19}. \url{https://doi.org/10.4171/owr/2023/19}. The supplied registry labels the reconstructed statement \texttt{corrected\_verified} and says it was checked against the official report; that verification has not been independently performed here.
\item Supplied literature source: Navarro--Ruhstorfer--Tiep--Vallejo, ``fields of values.'' \url{https://arxiv.org/abs/2304.12869}.
\item Supplied literature source: ``Oberwolfach Report 19/2023.'' \url{https://doi.org/10.4171/owr/2023/19}.
\item Supplied literature source: ``On a conjecture of Navarro and Tiep on character fields.'' \url{https://doi.org/10.1080/00927872.2025.2584406}.
\item Literature assessment dated exactly \texttt{2026-08-22}: supplied dated curation, not independently verified. It describes this continuation as the same fields-of-values Conjecture B, labels it partially solved, and states that one direction is established, the converse is reduced to quasisimple blocks, and no general proof is available. The background explicitly calls this dated literature triage derived from the supplied OWR curation and requests independent verification of time-sensitive status and cited claims.
\end{itemize}
\end{document}
